% ABSTRACT: one paragraph, at most 200 words, no citations or equations; define abbreviations.
% 2026-10-29 (user): the original opening (film bonded at the anneal and cooled while clamped) kept, manufacturing as the
% motivation; tilt senses, free energies and the process result follow. Previous versions: git history.
{\changedcolor
The band gap of a tin halide perovskite absorber is set by the film in the device, which bonds to its substrate at the
anneal and is cooled while clamped. How the film responds depends on its octahedral tilts, which differ between
compounds: caesium and formamidinium (FA) tin iodide tilt in-phase, methylammonium (MA) tin iodide antiphase. We
simulate \CsSnI{}, \MASnI{} and \FASnI{} with our machine-learned force field UrsaMLFF, fine-tuned on meta-GGA density
functional theory labels and about a hundred times faster than MACE models. On cooling, all three compounds
form their experimental tilt systems, with tetragonal tilt angles within \ang{0.3} of experiment. Free-energy calculations explain the opposite senses: antiphase tilting crosses the long axes of the upper
and lower faces of the cation cage, which frustrates FA, with an amine group at each end, but suits MA, which binds to
one face; \FASnI{} favours in-phase tilting by \SI{1.1}{\milli\hartree} and \MASnI{} antiphase tilting by
\SI{0.6}{\milli\hartree} per formula unit. Under clamping, the residual tension opens the gap of the
untilted hybrids elastically, by 8 to \SI{15}{meV}. In tilted \CsSnI{} it lowers the barrier of iodide migration by up
to \SI{35}{meV} per percent of strain, but barely shifts the balance between its black and yellow phases.
}
